%%%%%%%%%%%%
%
% $Autor: Müller, Hinrichs $
% $Datum: 2026-06-20 10:37:32Z $
% $Pfad: UltraschallradarAutomatisierungstechnik/Manuals/Benutzerhandbuch/Chapters/de/MainFunction.tex $
% $Version: 1 $
% !TeX spellcheck = en_GB/de_DE
% !TeX encoding = utf8
% !TeX root = manual 
% !TeX TXS-program:bibliography = txs:///biber
%
%%%%%%%%%%%%

\chapter{Die Hauptfunktionen}

\begin{figure}[h]
	\centering
	\begin{tikzpicture}[scale=1.6, transform shape, font=\sffamily]
	% Sensor
	\filldraw[fill=gray!30, draw=black, thick] (0,-1) rectangle (1,1);
	\node at (0.5, 0) [rotate=90] {\textbf{Sensor}};
	
	% Hindernis
	\filldraw[fill=red!30, draw=red!80!black, thick] (7,-1.5) rectangle (7.5,1.5);
	\node at (7.25, 0) [rotate=-90, text=red!80!black] {\textbf{Hindernis}};
	
	% Wellen Senden (Solid)
	\foreach \r in {1.5, 2, 2.5, 3} {
		\draw[thick, blue] (1,0) ++(45:\r) arc (45:-45:\r);
	}
	\draw[-Latex, blue] (3.5, 0.2) -- (5.5, 0.2) node[midway, above] {Sendesignal (Ping)};
	
	% Wellen Empfangen (Dashed)
	\foreach \r in {1, 1.5, 2} {
		\draw[thick, dashed, orange] (7,0) ++(135:\r) arc (135:225:\r);
	}
	\draw[-Latex, dashed, orange] (6, -0.2) -- (4, -0.2) node[midway, below] {Echo (Reflexion)};
\end{tikzpicture}
	\caption{Funktionsweise}
\end{figure}


Das Radarsystem arbeitet nach dem Prinzip des „Echo-Lot-Verfahrens“. Hierbei erfüllt das Gerät drei Kernaufgaben vollautomatisch:

\begin{enumerate}
	\item \textbf{Die Suchbewegung:} Die Dreheinheit schwenkt den Ultraschallsensor kontinuierlich in einem Halbkreis ($180^\circ$). Dadurch wird der gesamte Bereich vor dem Gerät lückenlos überwacht.
	\item \textbf{Die Messung:} Der Ultraschallsensor sendet für das menschliche Ohr unhörbare Schallimpulse aus. Treffen diese auf ein Hindernis, werden sie zurückgeworfen. Das Gerät misst die Zeit, bis das Echo eintrifft, und errechnet daraus zentimetergenau die Entfernung.
	\item \textbf{Die Visualisierung:} Diese Daten werden an Ihren Computer gesendet, wo sie in eine grafische Benutzeroberfläche übersetzt werden.
\end{enumerate}